\newpage
\section*{\S~5a Beförderungen}
\label{sec:befroerderungen}
\addcontentsline{toc}{section}{\S~5a Beförderungen}
\begin{enumerate}[label=(\arabic*)]
    \item Der größte Teil der Beförderung richtet sich nach der Dauer der aktiven Mitgliedschaft. Die genauen Voraussetzungen werden im Nachfolgenden aufgelistet. Als aktive Mitgliedschaft zählt die Zeit, die ein Mitglied nicht in der Reserve war.    
    \item Die Beförderungen innerhalb der Mannschaften, Unteroffiziere ohne Portepee und Unteroffiziere mit Portepee, mit Ausnahme der Dienstgrade Stabsfeldwebel und Oberstabsfeldwebel, werden durch die Leitung durchgeführt. Alle anderen Beförderungen werden durch die Mitgliederversammlung durchgeführt. Die Mitgliederversammlung kann durch eine relative Mehrheit ebenfalls Beförderungen innerhalb der Dienstgrade der Mannschaften, Unteroffiziere ohne Portepee und Unteroffiziere mit Portepee auf Vorschlag eines Mitgliedes vornehmen. Bei Beförderungen durch die Mitgliederversammlung innerhalb der Mannschaften, Unteroffiziere ohne Portepee und Unteroffiziere mit Portepee sind die unten genannten Voraussetzungen nicht nötig, die Voraussetzungen nach \hyperref[sec:dienstgradschluessel]{\S~5b} allerdings schon. Hierbei können auch passive Mitglieder befördert werden.    
    \item Bei Regelverstößen kann eine Aussetzung von Beförderung vorgenommen werden.    
    \item Dienstgrade können nicht übersprungen werden. Ausnahme hierbei sind: der Dienstgrad Korporal, welcher ab dem Dienstgrad Hauptgefreiter erreicht werden kann, der Dienstgrad Unteroffizier, welcher ab dem Dienstgrad Hauptgefreiter erreicht werden kann, der Dienstgrad Feldwebel, der ab dem Dienstgrad Hauptgefreiter erreicht werden kann, und die Dienstgrade Leutnant und Stabsfeldwebel, wenn diese durch die Wahl zu einem bestimmten Posten erreicht werden.   
    \item Es gibt kein Recht auf Beförderung. Beförderungen sind nur in Anwesenheit der zu befördernden Person durchführbar.
    \item Voraussetzungen:
    \begin{enumerate}[label=(\Alph*)]
        \item Voraussetzungen für Beförderungen innerhalb der Offiziere:
        \begin{center}
        \begin{tabular}{|p{4cm}|p{8cm}|}
        \hline
        \textbf{Dienstgrad} & \textbf{Voraussetzung} \\
        \hline
        Hauptmann & Dienstgrad Oberleutnant für mindestens 6 ununterbrochene Monate, mindestens eine vollständige absolvierte Dienstzeit als Unit-Leiter und Beförderung durch die Mitgliederversammlung \\
        \hline
        Oberleutnant & Dienstgrad Leutnant für mindestens 6 ununterbrochene Monate und Beförderung durch die Mitgliederversammlung \\
        \hline
        Leutnant & Wahl zum Unit-Leiter durch die Mitgliederversammlung, oder Dienstgrad Oberstabsfeldwebel für mindestens 6 ununterbrochene Monate, Unit-Leiter im Dienstgrad Hauptmann und Beförderung durch die Mitgliederversammlung \\
        \hline
        \end{tabular}
        \end{center}
	\newpage
        \item Voraussetzungen für Beförderungen innerhalb der Unteroffiziere mit Portepee (Feldwebel):
        \begin{center}
        \begin{tabular}{|p{4cm}|p{8cm}|}
        \hline
        \textbf{Dienstgrad} & \textbf{Voraussetzung} \\
        \hline
        Oberstabsfeldwebel & Dienstgrad Stabsfeldwebel für mindestens 6 ununterbrochene Monate und Beförderung durch die Mitgliederversammlung \\
        \hline
        Stabsfeldwebel & Wahl zum Spieß durch die Mitgliederversammlung \\
        \hline
        Hauptfeldwebel & Dienstgrad Oberfeldwebel für mindestens 6 ununterbrochene Monate oder Dienstgrad des Oberfeldwebels und aktive Mitgliedschaft für mindestens 24 Monate \\
        \hline
        Oberfeldwebel & Dienstgrad Feldwebel für mindestens 6 ununterbrochene Monate oder Dienstgrad Feldwebel und aktive Mitgliedschaft für mindestens 12 Monate \\
        \hline
        Feldwebel & Mindestens Dienstgrad Hauptgefreiter, am Tag der Feldwebelprüfung für mindestens 6 Monate aktive Mitgliedschaft und letzte passive Mitgliedschaft vor mindestens 6 Monaten, die Gruppenführer-Kampfrolle und bestandene Feldwebelprüfung \\
        \hline
        \end{tabular}
        \end{center}
        \item Voraussetzungen für Beförderungen innerhalb der Unteroffiziere ohne Portepee:
        \begin{center}
        \begin{tabular}{|p{4cm}|p{8cm}|}
        \hline
        \textbf{Dienstgrad} & \textbf{Voraussetzung} \\
        \hline
        Stabsunteroffizier & Dienstgrad Unteroffizier für mindestens 6 ununterbrochene Monate, oder besondere Leistungen als UvD \\
        \hline
        Unteroffizier & Mindestens Dienstgrad Hauptgefreiter und bestandene Unteroffiziersprüfung \\
        \hline
        \end{tabular}
        \end{center}
	\newpage
        \item Voraussetzungen für Beförderungen innerhalb der Mannschaften:
        \begin{center}
        \begin{tabular}{|p{4cm}|p{8cm}|}
        \hline
        \textbf{Dienstgrad} & \textbf{Voraussetzung} \\
        \hline
        Stabskorporal & Dienstgrad Korporal für mindestens 6 ununterbrochene Monate, Ausbilder für ingesamt drei Kampfrollen, wovon mindestens eine Kampfrolle eine Sonder-Kampfrolle ist und die Sch\"{u}tzen-Kampfrolle nicht mitgez\"{a}hlt wird, Rolle des Zeus/Missionsbauer und die Gruppenführer-Kampfrolle \\
        \hline
        Korporal & Mindestens Dienstgrad Hauptgefreiter, Ausbilder für mindestens zwei Kampfrollen inklusive der Bereitschaft bei administrativen Aufgaben zu unterstützen und die Kampfrolle Schütze auszubilden \\
        \hline
        Oberstabsgefreiter & Aktive Mitgliedschaft für mindestens 24 Monate oder 20 ununterbrochene Monate aktive Mitgliedschaft \\
        \hline
        Stabsgefreiter & Aktive Mitgliedschaft für mindestens 12 Monate oder 10 ununterbrochene Monate aktive Mitgliedschaft \\
        \hline
        Hauptgefreiter & Aktive Mitgliedschaft für mindestens 6 Monate \\
        \hline
        Obergefreiter & Aktive Mitgliedschaft für mindestens 3 Monate \\
        \hline
        Gefreiter & Kampfrolle Schütze und Teilnahme an drei Missionen als Jäger \\
        \hline
        Jäger & Mitgliedschaft in der 1st Aircat Division \\
        \hline
        \end{tabular}
        \end{center}
    \end{enumerate}
\end{enumerate}

\newpage
